% エントリ 13(2026-11-10): E12 を受けた論文の結論の整理
\section{[E13, 2026-11-10] E12 のあとの整理: 論文の結論はどう変わるか}\label{sec:e13}

E12 の訂正と指摘が、論文の主張(\texttt{paper/outline.md})にどう効くかをまとめる(ユーザー指示「研究ノートに書いといて」)。新しい計算はなく、数値は E9、E10、E12 のもの。すべて $\mu=0.71\times10^{-12}\eV$、主要項の $\gamma_\mathrm{BH}$、1 世代目の年齢 6 Gyr。

\subsection{変わらないこと}

\begin{itemize}
\item \textbf{窓の数値}: E12 の訂正 1(211 の臨界より下への降下)が効くのは $f\gtrsim10^{14}\GeV$ で、窓($f\sim10^{12}\GeV$)の数値は変わらない。
\item \textbf{論文の問い}: 「GWTC-5 の集団が、重い 1 世代目だけが回転を失う形を許すか」。どちらの結果でも論文として立つ構成も同じ。
\item \textbf{新しさ}: 自己相互作用を入れて集団の $(\mu,f)$ を走査した解析は、まだない。GW のスピンだけの上限は X 線連星の上限と同じ桁で、単独の事象からの制限としては強くない(E12)。新しさは集団解析そのものにある。
\end{itemize}

\subsection{変わったこと(窓が残る側に不利)}

\paragraph{GW241011 の主星の年齢}
E10 で「現実的」とした $T_2=5\times10^7\yr$ は $(50+50)\Msun$ の連星の遅れの最小値の目安で、この主星に近い $(10+10)\Msun$ では $3\times10^8\yr$ ほどになる。$3\times10^8\yr$ では、$f$ の上限は $4.6\times10^{11}\GeV$(0.72 で生まれる)--$7.6\times10^{11}\GeV$(0.78)、回転を失う 1 世代目は $17$--$22\Msun$ 以上($\chi_0=0.7$、$\delta=0.05$--$0.02$)になる(E12 の表~\ref{tab:e12age})。
\begin{itemize}
\item E10 で X 線連星を外して広がった分(GW だけで $11.7$--$14.5\Msun$ 以上、$\delta=0.18$--$0.05$)は、ほぼ失われる。
\item $\delta=0.05$ で比べると、年齢を直した GW だけの最小質量 $16.8$--$20.1\Msun$ は、E10 の GRS 1915+105 込みの現実的な場合の $19.3\Msun$ と同じ程度である。
\end{itemize}

\paragraph{判定の幅 $\delta$}
「1 世代目が臨界 $+\delta$ まで回転を失う」という判定は、下限が $f_\mathrm{low}\propto\delta^{-1/4}$ で $\delta$ に依存し、$\delta\to0$ では窓は閉じる(E12)。窓が残るかどうかは図からは言えず、集団解析で $\chi_\mathrm{SR}$ を測定と直接比べて決めるしかない。

\paragraph{1 世代目の自然スピン}
自然スピンが低い(Beta(0.69, 2.28))と、$10$--$20\Msun$ の 1 世代目の 6--8 割は 211 の臨界より下で生まれる(E12 の表~\ref{tab:e12zone})。窓が残る場合、積み上がりの信号は、重くて速く生まれた少数の 1 世代目から来ることになる。

\paragraph{逆向きの効果}
$\gamma_\mathrm{BH}$ の相対論的な補正(E7 の $C_\mathrm{BH}$)を入れると、GW241011 の主星からの上限は約 1.6 倍に上がり、1 世代目の最小質量は 1 $\Msun$ ほど下がる(E12)。窓は少し戻る。

\subsection{論文の書き方で変わること}

\begin{itemize}
\item 結果は年齢と組にして($\tau=\mu(f/\Mpl)^2T$ の形で、または年齢の事前分布の下で)示す。判定の幅 $\delta$ は使わない(DESIGN.md D9)。
\item 「2 世代目は較正された時計」という言い方はやめる。
\item 集団解析の前に、322 から育つ BH(E12 の訂正 3)を扱えるようにする。いまのままでは 1 世代目の標本の 1--2 割が模型の外になる。
\end{itemize}

\begin{result}
論文の結論そのものは、集団解析をするまで決まらない。ただ理論の側の見通しは、「現実的な仮定では窓はほぼ閉じており、残るとしても重い($\gtrsim17$--$22\Msun$)1 世代目だけが回転を失う狭い場合」のほうに傾いた。E8 の楽観的な帯 $f\approx(1$--$2.5)\times10^{12}\GeV$ は、2 世代目の年齢が $3\times10^7\yr$ 程度より短い場合(球状星団の遅れの目安より短く、AGN 円盤や一部の N 体計算の例にあたる)にしか成り立たない。
\end{result}
